\chapter{Experimental Methodology}
\label{chap:methodology}
\chapterrunning{experimental methodology}

This chapter describes the common experimental design used across the completed phases of the
project. It defines the evaluation framework carried from simple adversarial games to benchmark
routing and, finally, to bounded real-world vehicle routing. The framework is deliberately
conservative. Each phase relies on deterministic validators, held-out evaluation panels, repeated
runs, and phase-specific technical validity checks before any substantive interpretation is
accepted.

\section{Research Design Overview}
\label{sec:methods-overview}

The dissertation reports a sequence of completed experimental campaigns rather than a single
monolithic study. The sequence was chosen so that the search object, benchmark structure, and
validator pressure become progressively more demanding. Early phases test whether iterative code
evolution can alter behavior at all. Later phases ask whether the same loop can preserve useful
adaptations on standard optimization benchmarks with stronger baselines and clearer notions of
feasibility.

Table~\ref{tab:methods-phases} summarizes the completed evidence base retained in the final
dissertation.

\begin{table}[t]
\centering
\scriptsize
\setlength{\tabcolsep}{2pt}
\begin{tabularx}{\textwidth}{
  >{\raggedright\arraybackslash}p{0.72in}
  Y
  >{\raggedright\arraybackslash}p{0.62in}
  >{\raggedright\arraybackslash}p{0.72in}
  Y
}
\toprule
Phase & Search object & Replicates & Conditions & Primary endpoint \\
\midrule
1--4 (simple games) & full deterministic game-policy code & 10 transfer runs plus manual review & 3 transfer environments; 3 audited curriculum recipes & held-out win rate and score margin \\
5 (routing replay transfer) & constrained routing heuristic code & 20 paired offsets per family & 4 replay conditions across TSP, ATSP, and CVRP & held-out benchmark optimality gap \\
6 (replay mechanism study) & constrained TSP heuristic code & 20 paired offsets & 7 replay variants & held-out TSPLIB optimality gap \\
7 (operator discovery) & modular TSP operator code & 20 paired offsets & 7 conditions & held-out TSPLIB gap plus transplant/ablation validation \\
8 (adaptive portfolio) & TSP controller over a frozen heuristic portfolio & 20 paired offsets & 8 conditions & held-out TSPLIB gap and selector regret \\
9 (real-world solver evolution) & bounded CVRP solver code & 20 paired offsets & 4 conditions & held-out penalized gap with feasibility pressure \\
9 closeout (budget control) & bounded CVRP solver code under matched budgets & 20 paired offsets & 6 conditions & held-out penalized gap with explicit learned controls \\
10 (cross-family synthesis) & family-level result summaries & archived official artifacts & completed simple games, routing, and bounded CVRP plus 1 descriptive direct Codex baseline & normalized interpretation rather than pooled win counting \\
\bottomrule
\end{tabularx}
\caption[Completed empirical phases]{Completed empirical phases discussed in the final dissertation. A crossed model-tier study
was scoped but is treated here as future work rather than as part of the completed evidence base.}
\label{tab:methods-phases}
\end{table}

The archival phase record is finer-grained than the bundled presentation used in the dissertation.
In the project archive, phases 1, 2, 2b, 3, and 4 are preserved as distinct simple-game campaigns.
The final manuscript groups them into one simple-game evidence block because they share the same
environment family and lead to one retained transfer-and-novelty conclusion. The phase-9 closeout
is retained separately because it is not a new task family, but a scoped mechanism study on the
same bounded CVRP regime under matched learned controls. Phase 10 is retained as a synthesis stage
rather than as a new benchmark family. A separate phase-11 model-strength factorial study is
preserved in the research archive, but it is not admitted into the retained evidence base because
the final dissertation does not rely on that scoped follow-up campaign.

\begin{figure}[t]
\centering
\begin{tikzpicture}[
  scale=0.82,
  transform shape,
  box/.style={draw, rounded corners, align=center, minimum width=1.02in, minimum height=0.6in},
  line/.style={-Latex, thick}
]
\node[box] (games) at (0,0) {Simple\\games};
\node[box] (routing) at (2.75,0) {Routing\\transfer};
\node[box] (mechanisms) at (5.5,0) {Mechanism\\studies};
\node[box] (cvrp) at (8.25,0) {Real-world\\CVRP};
\node[box] (closeout) at (11.0,0) {CVRP\\closeout};
\node[box] (synth) at (13.75,0) {Cross-family\\synthesis};
\draw[line] (games) -- (routing);
\draw[line] (routing) -- (mechanisms);
\draw[line] (mechanisms) -- (cvrp);
\draw[line] (cvrp) -- (closeout);
\draw[line] (closeout) -- (synth);
\end{tikzpicture}
\caption[Progression of the experimental program]{Progression of the completed experimental program, from low-cost game adaptation to
routing, mechanism studies, higher-level search objects, bounded real-world CVRP, and the final
matched-budget closeout on the same CVRP regime. Portfolio control is included in the
mechanism-study stage.}
\label{fig:methods-progression}
\end{figure}

\section{Common Code-Evolution Loop}
\label{sec:methods-loop}

All phases instantiate the code-evolution loop introduced in
Section~\ref{sec:intro-formal-setup}. At each epoch, the model receives a bounded interface, the
current incumbent or scaffold, and structured feedback from prior evaluations. It then proposes a
deterministic Python artifact: a game policy, routing heuristic, modular operator, controller, or
whole CVRP solver. That proposal is validated, executed on the training panel, and either rejected
or admitted into the evolving archive.

Several constraints are shared across families.

\begin{enumerate}[leftmargin=2em]
\item Generated code is deterministic and auditable. Imports, hidden file or network access, and
arbitrary external solver calls are disallowed.
\item Task-specific validators separate syntactic success from scientific usefulness. A candidate
must execute and return a structurally valid output before it can be considered for acceptance.
\item Held-out evidence remains outside the inner loop. Training probes, replay probes, and archive
summaries may influence search, but the final claims are based on frozen end-of-run evaluation on
held-out panels.
\item Fallback or errored generations are logged diagnostically. They are not silently counted as
successful improvements.
\end{enumerate}

These constraints matter because they narrow the interpretation of success. A generated program is
useful only when it survives validation, remains executable across repeated evaluations, and
improves held-out performance under the benchmark-specific metric.

\section{Family-Specific Search Objects and Benchmarks}
\label{sec:methods-families}

The shared loop is implemented differently across families because the scientific question changes
with the object of search.

In the simple-game phases, the search object is full policy code for a two-player environment. The
environments emphasize transfer and behavioral change rather than exact optimality. Held-out win
rate and score margin are therefore the relevant outcomes.

In the routing replay phases, the search object is a constrained heuristic scaffold. For symmetric
TSP and ATSP, the evolved code modifies bounded construction, perturbation, and local-search
behavior over benchmark instances derived from TSPLIB \cite{reinelt1991,lin1973,helsgaun2000}. In
the initial CVRP extension, the same replay logic is transferred to a constrained multi-route
scaffold evaluated on CVRPLIB-style instances \cite{clarke1964,uchoa2017}. Relative or penalized
gap is the primary metric because it standardizes performance across instances of different scales.

Phase 7 changes the search object again, from whole-solvers to \emph{modular operators}. The aim
is not simply to find a better full heuristic, but to ask whether a reusable perturbation,
candidate-pruning, or restart component can survive transplant and ablation checks. Phase 8 then
switches from operator evolution to \emph{algorithm selection}: the model evolves a bounded,
interpretable controller over a frozen portfolio of classical TSP heuristics, which places the
phase directly within the algorithm-selection tradition introduced by Rice \cite{rice1976}. Phase
9 moves to bounded whole-solver evolution on real-world CVRPLIB X instances \cite{uchoa2017},
where feasibility, repair logic, and runtime variability matter as much as route cost. The phase-9
closeout then keeps the same bounded CVRP regime and validator, but replaces the original
comparison set with matched learned controls designed to distinguish replay-specific value from
extra candidate budget or direct one-repair synthesis alone.

\section{Replication, Aggregation, and Statistical Reporting}
\label{sec:methods-reporting}

The reporting discipline follows the conservative empirical guidance discussed in
Section~\ref{sec:lit-methodology}. The routing, operator, portfolio, and real-world CVRP phases
were all run with fixed paired seed offsets, usually twenty offsets shared across compared
conditions. This design supports paired deltas of the form in Equation~\eqref{eq:lit-paired-delta}
and reduces the chance that one condition appears stronger only because it benefited from an easier
stochastic draw.

The aggregate reports emphasize four summary types.

\begin{enumerate}[leftmargin=2em]
\item Condition means on the primary held-out endpoint.
\item Paired condition deltas with 95\% bootstrap intervals.
\item Technical-validity summaries, including fallback counts, generation errors, timeout checks,
and acceptance counts.
\item Diagnostic metrics such as code novelty, adaptation efficiency, archive hardness, or runtime
inflation when those metrics are relevant to the specific phase.
\end{enumerate}

The dissertation does not rely on best-run anecdotes. When a phase produces a strong positive
result, that result is reported together with the paired uncertainty and with the corresponding
validity checks. When a phase produces a null or boundary result, that outcome is treated as a
finding rather than as a failed attempt to be hidden.

\section{Qualitative Audit and Validity Safeguards}
\label{sec:methods-validity}

Not every important question can be answered by scalar metrics alone. Two additional safeguards are
therefore used.

First, the simple-game transfer work includes a manual novelty adjudication pass. The largest code
novelty spikes are reviewed and classified as functional adaptation, mixed or local adjustment, or
churn/failed adaptation. This guards against the common mistake of equating large lexical change
with strategically meaningful change.

Second, the later routing phases use phase-specific technical validity audits before interpretation.
Examples include checking that all paired runs and all conditions are present, verifying that
accepted artifacts did not depend on fallback behavior, and confirming that operator or solver
validation pipelines executed as intended. The phase-7 operator campaign, for instance, was deemed
scientifically usable only after confirming zero modular generation fallbacks and zero
materialization fallbacks in the official run set. Phase 9 was likewise interpreted only after
checking full held-out feasibility, zero accepted fallback generations, and the absence of hidden
solver-worker timeout contamination in the completed official runs. The phase-9 closeout adds one
more safeguard to the same regime: replay is credited only after comparison against matched learned
controls that remove the simpler explanations of extra search budget or fresh direct synthesis.

\section{Scope of the Retained Evidence}
\label{sec:methods-scope}

The final manuscript is intentionally restricted to completed phases whose evidence is coherent and
whose reporting pipeline is already established. One descriptive exception is retained: a direct
single-shot Codex-authored CVRP solver evaluated through the same phase-9 validator. That artifact
is informative as a sanity baseline for interactive coding assistance, but it is not pooled with
the replicated stochastic campaigns because it is neither randomized nor generated under the same
frozen API-loop protocol.

This scope choice reflects the dissertation's broader methodological position. The objective is not
to maximize the number of claimed wins, but to produce a defensible account of when the code
evolution loop adapts, when it churns, and how those patterns change as validators and baselines
become harder.

For reproducibility, the retained claims in this dissertation are tied to official archived run
artifacts and code snapshots preserved in the project version-control record. Where stable
versioned snapshots exist, such as phase-completion tags, those archived states are the intended
reference points for the retained evidence rather than a moving development branch.
